\section{Discussion}
\label{sec:discussion}

At the main matching margin, Kairo's strongest overall profile matched more annotated changes than the strongest \softwarepm{} profile. The \softwarepm{} profile achieved a better balance between precision and recall at every margin tested. Additional alarms give an analyst more intervals to investigate and increase the effort required to monitor a process. Both profiles also raised alarms on logs without injected changes. Those alarms may reflect ordinary fluctuations or other changes outside the benchmark annotations, and their operational importance would need separate investigation.

The examples help explain when a state representation is useful. A resource profile surfaced a reassignment, and an inter-case profile surfaced a gradual increase in waiting time. Activity sequences alone provide little direct evidence about either change. Where the injected change did affect activity sequences, methods that compare sequences or activity relations performed well. In particular, Kairo and version clustering gave the same aggregate results on the selected Padella control-flow logs. The case-composition experiments illustrate another part of the picture. Once the comparison methods were given the relevant case attributes explicitly, they located the region changes cleanly. The choice of measurements can matter as much as the choice of the detector.

Separate state descriptions give analysts different views of the same episode. An analyst can inspect case progress, resource activity, and overall load around the same time and compare the measurements behind the states that became more or less common. Co-occurring signals can guide that investigation. For instance, a reassignment may change workload and waiting time together. The assessment compared the timing of alarms with annotations. Whether the interface helps analysts explain a change, or trace a causal link between perspectives, was outside that evaluation.

The benchmark was retrospective. Kairo learned the states and the relative alarm threshold from the full log, and the strongest profiles were identified on the logs used for scoring. Its performance therefore describes retrospective localization on the selected benchmark. The alarm threshold marks large changes relative to the rest of the same log. A stable process can still produce peaks, and a log with several large changes can provide a demanding reference for smaller ones. The log families were chosen for coverage of different behaviors and change types. The comparison cannot establish how often the same results would occur in operational logs or with prospective monitoring.
